\section*{Appendix}

\subsection*{A. Method Changes Relative to the Earlier Protocol}
This appendix records superseded definitions so that numbers reported in earlier
versions of this work can be located and not confused with the present ones. None of
the quantities below appear in the main text.

\paragraph{A.1 Event corpus.}
The earlier corpus was assembled by a geometric hazard predicate followed by a check
that the human braked. Its corridor criterion admitted events in which the entity was
already visible in the designated ``clean'' frame, so the contrast measured a closer
approach rather than an appearance. It is fully replaced by the brake-first
construction of \cref{sec:mining}.

\paragraph{A.2 Readout.}
The earlier readout was the policy's native \texttt{commanded\_speed}, covering
$0.25$--$0.75$\,s of a $2.5$--$4.0$\,s plan. On identical events the ground-truth
hazard-direction magnitude reads $0.332$ under that window and $1.232$ over the full
horizon, a factor of $3.7$. \texttt{commanded\_speed} is retained only as a
sensitivity readout.

\paragraph{A.3 Verdict system.}
The earlier three-state PASS / FAIL / indeterminate adjudication is replaced by the
continuous response ratio of \cref{eq:ratio} with an interval, which does not require
a threshold on a quantity whose scale differs across policies.

\subsection*{A.4 Deployment metrics, written out}
The official aggregate combines four multiplicative terms with a weighted sum of four
others:
\begin{equation}
\mathrm{PDM} \;=\;
\Big(\!\!\prod_{k\in\mathcal{M}}\! m_k\Big)\cdot
\frac{\sum_{j\in\mathcal{W}} w_j\,s_j}{\sum_{j\in\mathcal{W}} w_j},
\label{eq:pdm}
\end{equation}
with $\mathcal{M}=\{$no at-fault collision, drivable area, driving direction, traffic
light$\}$ and $\mathcal{W}=\{$progress, TTC, lane keeping, comfort$\}$. The progress
term is normalised \emph{within the proposal batch},
\begin{equation}
s_{\mathrm{progress}}^{(i)}=\mathrm{clip}\!\left(
\frac{p_i}{\max_{i'} \big(p_{i'}\!\cdot\!\prod_k m_k^{(i')}\big)},\,0,\,1\right),
\label{eq:progress}
\end{equation}
which is identically $1$ when a single proposal is scored alone --- the degeneracy
reported in \cref{sec:validity}. We therefore score all candidates as one batch, which
makes progress \emph{relative} to the batch and not comparable with published absolute
figures.

Our own comfort readouts are taken from the planned path resampled at
$\delta=0.5$\,s, with $u_k$ the speed sequence of \cref{eq:onset}:
\begin{equation}
a_k=\frac{u_{k+1}-u_k}{\delta},\quad
j_k=\frac{a_{k+1}-a_k}{\delta},\quad
J=\sqrt{\tfrac{1}{n}\textstyle\sum_k j_k^{2}} .
\label{eq:jerk}
\end{equation}

\subsection*{B. Corrections Made During This Study}
\paragraph{B.1 Velocity frame (deployment corpus).}
Object velocities in the deployment logs are expressed in the ego frame, not the
global frame; an earlier version applied a further rotation, which made every object
appear stationary to the attribution of \cref{eq:areq} and inflated the lead-vehicle
pool by $3.7\times$. The benchmark corpus derives velocity by differencing global
positions and was unaffected (verified: pool sizes unchanged).

\paragraph{B.2 Turn gate.}
A heading-change gate estimated from path tangents produced spurious $\sim180^{\circ}$
turns whenever the ego came to a full stop, because the trailing path points coincide.
Replacing it with the unwrapped ego-pose yaw plus a $5$\,m displacement gate removed
$33$ of $38$ events from the affected pool.

\paragraph{B.3 Scene identifiers.}
Scene names are not unique across dataset splits: $65.8\%$ of the deployment test
split's scene names also occur in the trainval split, referring to different logs.
All merged pools are keyed by (split, scene, frame).

\paragraph{B.4 Sufficiency of the patched set.}
See \cref{sec:validity}.
